% One Britannica scan as extracted by Nemotron Parse 3.5. Filled by src/nemotron_challenge/stages/latex.py.
% The TITLE, BODY and SOURCE markers (each wrapped in two percent signs) are replaced. Compile with XeTeX (tectonic).
\documentclass[10pt]{article}
\usepackage[a4paper,margin=18mm]{geometry}
\usepackage{fontspec}
\usepackage{multicol}
\usepackage{graphicx}
\usepackage[hidelinks]{hyperref}
\setlength{\parindent}{0pt}
\setlength{\parskip}{4pt}
\setlength{\columnsep}{8mm}
\begin{document}
\section*{Historical source: scan row 1678}
\begin{multicols}{2}
\subsection*{\textit{continued from the previous page}}
{\small\itshape [Margin: Blfinger.]}\par
Billinger. Edward, had already had occasion to appreciate his talents, and put them to use. At the time when he carried on the war in Serie, he maintained a regular correspondence with Billinger who had long been known as an able Engineer, and had, indeed, made some improvements of value in the received system of Portification. After his return to Tubingen, he had frequent conversations with the Professor, on different subjects of administration and appointed him, in 1735, Privy-Comseller. This nomination was not a simple honorary title. Billingers saw himself raised at once to a power almost unlimited. He revisited some times pro motion, which he did not think himself qualified to sustain. In accepting office, his first case was to require all the knowledge necessary to the discharge of its duties. He employed almost two years in assiduous laborto to interact himself thoroughly in the statistics of the country, in considering its political situation, its constitution, its interests, and became, at the end of all this study, one of the most enlightened ministere that his country had yet proved. Billinger was placed in a situation too elevated not to excite jedowny and harred. He left is, and waited to get the ministry; but the court refused to receive his resignation, soon after the tender of which, the Duke died. Billinger experience from his successor of the consideration and all the friendship which he had experienced in the beginning of his career. Received into a confidence without bounds, he had the power to realize, without obstacle, those plans of administration with which the most enlightened particulars had inspired him. Wurtemburg still feels the happy influence of his ministry. Commerce, public interaction, agriculture, were protected and amorliorated by his cases. The culture of the true, of so much importance in this country, was one of the principal objects of his attention. We ought not to forget that he was the original author of that strict union which has long united Wurtembourg and Prussia, and of the importantes to which the hereditary Prince of Wurtemburg was raised at the court of Belfin. In 1737, the Duke nominated him President of his consolony, and Secretary of the grant Order of the Chace. He was also Counter or the University of Tubingen, and member of the Royal Academy of Berlin. All his time was consecrated to some serious occupation, with the exception of one hour in the evening, which he employed in making and receiving visits. His greatest employment was in cultivating his garden. A warm and strong friend, he gave many proofs of gratitude to those protectors, who had previously assisted him in his studies. He has been reproached with being insecible; but in spite of some slight blemishes, the memory of Blfinger will be always due to his countrysmen, and honored by all Germans. Wurtemburg reckons him among the greatest men which she has produced, and proposes him as a model to the statement and her man of letters. He was more married, and left no issue. He died at Settargar the 18th of February 1750. His works, besides various Papers published in the Memoirs of the St Petersburg and Paris Academies of Science, are : 1. \emph{Disputatio de Harmoniá prestabilità}, Tubingen, 1723, in 40c. 2. \emph{De Harmonia Animi et Corporis Humani mazi mè prestabilità Commentato Hypothetica}. Francfort on the Rain, 722, in 8vo. This work was inserted in the Exrangatory Index in Rome in 1734. 3. \emph{De Origins et Permissione Mali, precipue Morrulla, Commentato Philosophica}, ibid. 1724, in 8vo. 4. \emph{Specimen Dietrich Veteran Sinatur Moralis et Politica}, francf. 1724, in 40c. 5. \emph{Discretato Historiocetapria da Speculo Arionatie}, Tubingen, 1725, in 40c. 6. \_Dilucidations Philosophica de Deo, Anima Himmoni, Monato, et Generatious Revarum Affectioni-Ion\_, Ibid. 1725, in 40c. 7. \emph{Billingert et Holmanni Epistola de Harmonia Prestabilità}, 1728, in 40c. 8. \emph{Disputate de Mararó et Legius Studii in Theologiä Thetici}, ibid. 1731, in 40c. 9. \emph{Disputate de Collo De Rationali}, ibid. 1731, 10. \emph{Victor Breves in Losa, Spinose Melodam Explicanti Scriptures}, Tub. 1732, in 40c. 11. \emph{De Mysteris Gratitura Piala Generatin Spectatsi Sorneo, Nectitatus} 1733, Tubingen, 1728, in 40c. 12. \emph{Elementa Physics}, Leipzig, 1742, in 8vro. 13. \emph{La Ciudelle Coupé}, Leipzig, 1756, in 40a. See \emph{Biographie Universelle}, Tom. IV. (z.)
\end{multicols}
\section*{Source record}
\begin{itemize}
\item Page ID: \texttt{britannica:309f765b1ecc83e0a7f7762e44b4a927a7555443:train:1678}
\item Scan: \texttt{page-1678-75b494c1b4fb.jpg}, 600 x 817 px
\item Image SHA-256: \texttt{\small 75b494c1b4fb241bd9dcaba505dfb1732b3c54ff9f843f0d4a7c1c31110a6ec8}
\item Extraction: \texttt{nvidia/nemotron-3.5-parse}, raw reply \texttt{artifacts/extract/raw/1678.json}
\item Text is Parse OCR output, normalized (long s as s, line-end hyphens joined). It is not corrected by hand.
\end{itemize}
\end{document}
